\chapter{Logica del prim'ordine}

La logica proposizionale si concentra essenzialmente sulle proprietà degli operatori logici, ed è detta proposizionale perché sono considerati atmomiche le proposizioni. Questo si riflette nella traduzione dal linguaggio naturale, in cui la separazione tra proposizioni avveviva tramite congiunzioni, disgiunzioni, negazioni, implicazioni, ecc.

Tale granularità però non è sufficiente per esprimere relazioni tra elementi del discorso interni alle proposizioni. Consideriamo nuovamente infatti l'esempio del sillogismo aristotelico \q{Ogni uomo è mortale, Socrate è un uomo, dunque Socrate è mortale}. Una traduzione in logica proposizionale potrebbe essere \((p\land q)\to r\) poiché le tre frasi internamente non hanno una struttura proposizionale. Chiaramente il fatto che tali frasi siano considerate valide non può dipendere dalla loro struttura proposizionale, poichè la struttura proposizionale non è sufficiente per esprimere la validità del sillogismo. Ciò che intuitivamente collega le frasi e rende sensata la validità sono elemnti interni alle proposizioni, come il fatto che \q{Socrate} è un \q{uomo} e che ogni \q{uomo} è \q{mortale}. Di conseguenza per esprimere la validità di tale sillogismo è necessario andare ad analizzare una struttura delle frasi più granulare che nella logica proposizionale viene incporporata negli atomi.

In questo caso la relazione che ci viene nascosta è il fatto che \q{uomo} è una proprietà degli ogetti, che può essere vera o falsa, stessa cosa per \q{mortale}. Inoltre \q{Socrate} è un elemento specifico del nostro universo. Quello che quindi ci dicono queste frasi è che un ogetto del nostro dominio che è \q{uomo} è \q{mortale}, e che in particolare l'ogetto \q{Socrate} è un \q{uomo} e dunque se ne deduce che \q{Socrate} è \q{mortale}. Il salto concettuale quindi, da logica proposizionale al prim'ordine, è l'arricchimento del concetto intuitivo di situazione da un associazione tra frasi e valori di verità che mi indicano se sono accaduti o meno, all' associazione tra gli ogetti del discorso e le proprietà che essi hannno.

Senza entrare subito nel dettaglio della semantica, per esprimere tale granularità si avrà necessariamente di nomi di elementi di una situazione, quali possono essere anche finiti in un universo di ogetti infiniti. I nomi non sono, come nella logica proposizionale, ogetti sintattici che vengono interpretati in valori di verità, ma sono oggetti sintattici che vengono interpretati in oggetti del dominio. Dunque, anche il concettto di interpretazione cambia, diventando più complesso.

Inoltre, per parlare di tali ogetti sarà necessario parlare di proprietà di tali ogetti, o di relazioni tra essi nel caso siano coinvolti più ogetti nel discorso. Per proprietà e relazioni, il concetto che esse esprimono è molto vicino al concetto matematico di proprietà. Ad esempio, per i numeri una proprietà è quella di essere pari. Ovviamente da sola non ha molto senso, ma applicata ad uno specifico ogetto del dominio come ad esempio \(5\) è interpretabile con un valore di verità, ovvero vero se l'ogetto ha quella proprietà, falso altrimenti. Ovviamente tale assegnamento dipende dall'interpretazione che diamo alla stringa \q{pari}, che dal punto di vista logico non è altro che un simbolo sintattico (o un nome di proprietà) a cui possiamo dare una semantica arbitraria. Senza un modello, \q{pari} è solo un simbolo, e non ha alcun significato.

Infine, è usuale nel linguaggio parlare di ogetti identificandoli indirettamente, tramite una dipendenza con altri ogetti. Di conseguenza sarà necessario avere modi di costruire espressioni che si riferiscono ad ogetti a partire da altri ogetti. Considerando ad esempio l'espresisone \q{Il cane di Socrate aveva 3 gambe}, la locuzione \q{il cane di Socrate} identifica un ogetto del dominio come farebbe un nome, ma senza nominarlo esplicitamente, ma identificandolo in funzione di un altro ogetto. Dunque possiamo rappresentarla sintatticamente con una funzione \(f\) che prende in input un ogetto del dominio, e restituisce un altro ogetto del dominio. Per quanto questo possa sembrare molto simile ad una proprietà, è importante notare che una proprietà è interpretata in un volore di verità, mentre una funzione in un ogetto del dominio. 

La logica del prim'ordine permette dunque di parlare di ogetti di un dominio direttamente (nomi detti costanti) o indirettamente a partire di altri ogetti (funzioni), e di parlare di proprietà di tali ogetti o relazioni tra essi (predicati). Questo tipo di costruizione di nomi ci permette di creare un insieme infinito di nomi a partire da un insieme finito di costanti e funzioni.\footnote{Per avere un insieme infinito di termini sono sufficienti una costante e una funzione unaria. Interpretando la costante con il numero \(0\) e la funzione come la funzione successore, è possibile infatti rappresentare tutto \(\N\).} L'insieme di questi ogetti è detto insieme dei \keyword{termini}. Tali termini possono essere usati come argomenti di proprietà, relazioni o altre funzioni creando nuovi termini.